\section*{Capítulo \ref{ch:b2:mass-spec-atomic} --- Espectrometria de massas e espectroscopia atômica}

\begin{solution}{exo:b2:mass-spec-atomic:1}
78; $400/2 = 200$; $194 + 1 = 195$, $[\mathrm M + \mathrm H]^+$.
\end{solution}

\begin{solution}{exo:b2:mass-spec-atomic:2}
Massa ímpar: um número ímpar de átomos de nitrogênio, um só para uma molécula tão pequena. \ce{C3H7NO}
(propanamida, 73) ou \ce{C4H11N} (butan-1-amina, 73).
\end{solution}

\begin{solution}{exo:b2:mass-spec-atomic:3}
Quase iguais: um bromo. 3 : 1: um cloro.
\end{solution}

\begin{solution}{exo:b2:mass-spec-atomic:4}
Amarelo: sódio, \qty{589.0}{nm} e \qty{589.6}{nm}. Carmim: lítio, \qty{670.8}{nm}.
% ledger: asd:NaD2, asd:NaD1, asd:Li671
\end{solution}

\begin{solution}{exo:b2:mass-spec-atomic:5}
$n_C \approx 6.7/1.11 \approx 6$.
\end{solution}

\begin{solution}{exo:b2:mass-spec-atomic:6}
\ce{CO}: 27.9949; \ce{N2}: 28.0061; \ce{C2H4}: 28.0313 u. Poderes de resolução: \ce{CO}/\ce{N2} $28/0.0112
\approx 2.5 \times 10^3$; \ce{N2}/\ce{C2H4} $28/0.0252 \approx 1.1 \times 10^3$; \ce{CO}/\ce{C2H4}
$28/0.0364 \approx 7.7 \times 10^2$.
% ledger: mass:C12, mass:O16, mass:N14, mass:H1
\end{solution}

\begin{solution}{exo:b2:mass-spec-atomic:7}
86: o \omterm{def:b2:mass-spec-atomic:molecular-ion}{íon molecular} de \ce{C5H10O}. 57: \ce{C2H5CO+}, um \omterm{def:b2:aromatic-substitution:friedel-crafts}{íon acílio} vindo de uma clivagem $\alpha$
(perda de um radical etila, 29). 29: \ce{C2H5+}.
% ledger: ms:pentan-3-one
\end{solution}

\begin{solution}{exo:b2:mass-spec-atomic:8}
86: $\mathrm M^{+\bullet}$. 71: $\mathrm M - \ce{CH3}$, o acílio \ce{C3H7CO+}. 43: \ce{CH3CO+}
(perda do radical propila). 58: \omterm{def:b2:mass-spec-atomic:fragmentation}{rearranjo de McLafferty}, por meio de um hidrogênio do carbono
$\gamma$ da cadeia propila, com perda de eteno. Na pentan-3-ona os grupos etila não têm carbono
$\gamma$: nenhum íon de McLafferty; seu pequeno pico em 58 é o pico isotópico \ce{^{13}C} do íon em 57
(três carbonos, cerca de 3.3~\%).
% ledger: ms:pentan-2-one, ms:pentan-3-one
\end{solution}

\begin{solution}{exo:b2:mass-spec-atomic:9}
Uma reta que passa pela origem tem a inclinação $\sum c_iA_i/\sum c_i^2 = 1.553/30.0 \approx \qty{0.0518}{L/mg}$. Amostra:
$0.128/0.0518 \approx \qty{2.47}{mg/L}$ de cobre.
\end{solution}

\begin{solution}{exo:b2:mass-spec-atomic:10}
$t = L\sqrt{m/(2eU)} = 1.00\sqrt{1000 \times 1.6605 \times 10^{-27}/(2 \times 1.6022 \times 10^{-19}
\times 2.00 \times 10^4)} \approx \qty{16.1}{\micro s}$. Para 1001 u, $t$ cresce pelo fator
$\sqrt{1.001} \approx 1 + 0.0005$: $\Delta t \approx \qty{8.0}{ns}$.
% ledger: const:u, const:e
\end{solution}

\begin{solution}{exo:b2:mass-spec-atomic:11}
Dois cloros: $(0.758 + 0.242x)^2$ dá $0.575$, $0.367$, $0.0586$: $m/z$ 84, 86 e 88 na razão
100 : 64 : 10.
\end{solution}

\begin{solution}{exo:b2:mass-spec-atomic:12}
\ce{^{40}Ar^{16}O+}: $39.9624 + 15.9949 = 55.9573$; \ce{^{56}Fe}: 55.9349; $R = 56/0.0224 \approx 2.5 \times
10^3$. Ou então meça outro isótopo do ferro, o \ce{^{57}Fe}, ou elimine o \ce{ArO+} em uma célula de
colisão antes do analisador.
% ledger: mass:Ar40, mass:Fe56, mass:O16
\end{solution}

\begin{solution}{pb:b2:mass-spec-atomic:1}
\textbf{1.} 156, 158 e 160 (com 157 e 159): a molécula existe em várias formas isotópicas.
\textbf{2.} Três picos separados de duas unidades, o do meio o maior: um cloro e um bromo.
\textbf{3.} 156.
\textbf{4.} Massa par: nenhum nitrogênio, ou um número par.
\textbf{5.} $156 - 79 - 35 = 42$: \ce{C3H6}.
\textbf{6.} \ce{C3H6BrCl}; grau de insaturação $(2 \times 3 + 2 - 6 - 2)/2 = 0$: nenhum anel, nenhuma
ligação dupla.
\textbf{7.} $0.5/14.5 \approx 3.4~\%$, $3.4/1.11 \approx 3$ carbonos.
\textbf{8.} O pico em 157 é dado com 0.1 de precisão em uma escala na qual ele vale 0.5: uma incerteza de 20~\%.
\textbf{9.} \ce{C3H6^{81}Br^{35}Cl} e \ce{C3H6^{79}Br^{37}Cl}.
\textbf{10.} $3 \times 12 + 6 \times 1.007825 + 78.918338 + 34.968853 \approx 155.9341$.
\textbf{11.} $156/(156.151 - 155.934) \approx 7.2 \times 10^2$.
\textbf{12.} Sim: dois hidrogênios a menos, 154.
\textbf{13.} Um átomo de bromo (79 ou 81) é perdido: \ce{C3H6Cl+}, 77 com \ce{^{35}Cl}, 79 com \ce{^{37}Cl},
3 : 1.
\textbf{14.} Perda de HBr: o cátion radical $\mathrm{C_3H_5Cl^{+\bullet}}$, de massa par.
\textbf{15.} \ce{C3H5+}, o cátion alila, após a perda dos dois halogênios (Br, depois HCl).
\textbf{16.} \ce{CH2Cl+} (49 e 51, 3 : 1).
\textbf{17.} \ce{CH2Br+} (93 e 95) e \ce{C2H4Br+} (107 e 109), cada um 1 : 1: bromo em uma extremidade
da cadeia, o segundo vindo da perda de \ce{CH2Cl}.
\textbf{18.} A ligação \ce{C-Br} é mais fraca que a ligação \ce{C-Cl}, e \ce{Br} é o melhor radical de
saída.
\textbf{19.} Não: nenhum pico em 127--131; \ce{CH2Cl+} e \ce{CH2Br+} mostram cada halogênio em um
\ce{CH2} terminal.
\textbf{20.} \ce{BrCH2CH2CH2Cl}, 1-bromo-3-cloropropano.
\textbf{21.} 156/158/160: M; 77/79: M $-$ Br; 76: M $-$ HBr; 107/109: M $-$ \ce{CH2Cl};
93/95: \ce{CH2Br+}; 49/51: \ce{CH2Cl+}; 41: \ce{C3H5+}; 39, 27: \ce{C3H3+}, \ce{C2H3+}.
\textbf{22.} Não: não há grupo carbonila.
\textbf{23.} 1-bromo-2-cloropropano, \ce{CH3CHClCH2Br}: a perda de \ce{CH2Br} dá \ce{CH3CHCl+} em 63/65.
\textbf{24.} \ce{C3H6^{81}Br^{37}Cl}.
\textbf{25.} $M$: $0.758 \times 0.5069 = 0.384$; $M+2$: $0.758 \times 0.4931 + 0.242 \times 0.5069 =
0.496$; $M+4$: $0.242 \times 0.4931 = 0.119$. Razão $\boldsymbol{\approx 3.2 : 4.2 : 1}$, cerca de 3 : 4 :
1; o espectro dá $14.5 : 18.6 : 4.4 = 3.3 : 4.2 : 1$.
\end{solution}
